% Introduction. Numbers via macros; claims follow evidence/discrepancies.md.

\section{Introduction}
\label{sec:intro}

Weakly supervised video anomaly detection learns to localize abnormal events in time from video-level labels alone. Current methods start from frozen snippet features, most often from CLIP~\cite{radford2021clip}, and train a light temporal model on top with a multiple-instance objective~\cite{sultani2018real}. VadCLIP~\cite{wu2024vadclip} extends this design with a second output that aligns snippets with learnable class prompts, and it is a common base for recent work~\cite{yang2024tpwng,wu2024stprompt}.

Two design choices in this pipeline limit what the model can learn. The first concerns the input. Before any learning, each training video is resampled to a fixed length, and long videos are mean-pooled over uniform segments. On the training sets we use, this affects \ucfTrainAvgFrac\% of UCF-Crime and \xdTrainAvgFrac\% of XD-Violence items. Changes between snippets inside a segment do not survive pooling, and differencing the pooled sequence only measures change between segment means. The second concerns the objective. VadCLIP's semantic output is computed downstream of its classification output, and both read the same representation. Adding a semantic multiple-instance loss to the whole network therefore also moves the classification output. In a single-seed run, doing so raised XD-Violence AP by \fullGXdApTwo{} points and changed UCF-Crime AUC by \fullGUcfAucOne{} points.

We address the two choices in two training stages. Stage~1 computes the absolute difference between consecutive raw snippets before pooling, aggregates it with the same segment boundaries as the appearance features, and injects it through a zero-initialized low-rank residual; it also mixes VadCLIP's similarity graph with a sparse $k$-nearest-neighbour graph. Stage~2 optimizes the semantic multiple-instance loss with respect to VadCLIP's text-side adapter only. That adapter lies downstream of the classification output, which is therefore unchanged by construction.

We evaluate under a protocol in which neither training nor checkpoint selection uses the test set, retrain VadCLIP under the same protocol, and report paired differences over three seeds. Our method improves UCF-Crime AUC by \ucfGainMean{} points and XD-Violence AP by \xdGainMean{} points, with the same sign on every seed. Computing the same variation after pooling instead falls below the baseline on every seed, and an equal-budget control without the semantic objective shows that the Stage~2 gain is not an effect of additional training. Attribution experiments locate the UCF-Crime gain in Stage~1 and the XD-Violence gain in Stage~2.

Our contributions are as follows.
\begin{itemize}
  \item We show that computing temporal variation before fixed-length pooling, rather than after it, changes what a CLIP-based detector learns from long training videos, and we inject it through a residual that adds \pvrParams{} parameters.
  \item We restrict semantic post-training to the part of VadCLIP downstream of its classification output, which keeps that output unchanged while improving the semantic output, without adding parameters.
  \item We evaluate both stages in a matched multi-seed comparison with an equal-budget control and component attribution, under a protocol in which neither training nor checkpoint selection uses the test set.
\end{itemize}
